% -----------------------------------------------------------------------------
% Section 6: Discussion
% Paper: "Silent Decisions Are Type Errors: Enforcing AI Accountability
%         via Linear Resource Types"
% Venue: IEEE Computer — Special Issue on AI Governance
% Deadline: 13 April 2026
% Authors: Daniel Lau Pereira Soares
% Rev 2: 24 March 2026 — four reviewer-hardening revisions:
%   Gap 1: AttestableContext trait for semantic truthfulness
%   Gap 2: FFI gateway + BTV sidecar for polyglot pipelines
%   Gap 3: ComplianceAuthority factory replaces public constructor
%   Gap 4: CAL demoted to design trade-off; CAP analogy removed
% Rev 3: 29 March 2026 — security audit patches applied:
%   Patch 1.1: Trust-assumption alignment with the log authority (BTV_HMAC_KEY)
%   Patch 1.2: ComplianceAuthority factory (implemented, not proposed)
%   Patch 1.3: Sixth future-work item — formal composition theorem
% Rev 4: 14 September 2026 — IEEE Computer submission hygiene:
%   E1: economics kept qualitative (no crossover volume, no companion cite)
%   R2.1: CAL compressed to one paragraph (design trade-off, not theorem)
%   E3: FFI threat model promoted into Section 4 (sec:ffiboundary); L4 points there
%   Sibling-manuscript citations removed; refs.bib orphan entries pruned
% -----------------------------------------------------------------------------

\section{Discussion}
\label{sec:discussion}

The Constitutional Enclosure Theorem establishes a strong structural
guarantee within a precisely bounded scope.
This section examines the boundaries of that scope honestly,
identifies limitations of the current reference implementation,
and proposes concrete extensions that preserve the core invariant
while closing each gap.

% ----------------------------------------------------------------
\subsection{Scope and Boundaries of the Guarantee}
\label{sec:scope}

\paragraph{Unsafe Rust.}
The \texttt{unsafe} keyword grants the programmer the ability to
bypass borrow-checker and visibility rules, including the private
field restrictions that enforce the encapsulation boundary
(Section~\ref{sec:encapsulation}).
Production deployments should enforce a zero-\texttt{unsafe} policy
over the transitive dependency graph, auditable via
\texttt{cargo geiger}~\cite{cargogeiger}.
This is an operational policy decision, not a theoretical barrier:
the tooling exists and its enforcement requires no changes to the
BTV crate.

\paragraph{Semantic truthfulness and the AttestableContext extension.}
The theorem guarantees that \emph{some} byte sequence was hashed
and bound to a verdict---not that the sequence faithfully represents
the context presented to the AI system at inference time.
A caller could pass \texttt{b"placeholder"} to
\texttt{EvidenceToken::new} and obtain a structurally valid token.

We propose a concrete extension that narrows this gap without
violating the core invariant.
Introduce a sealed \texttt{AttestableContext} trait:

\begin{lstlisting}[language=Rust,caption={AttestableContext trait
  (proposed extension)}]
pub trait AttestableContext: Sealed {
    fn canonical_bytes(&self) -> &[u8];
    fn schema_version(&self)  -> u16;
    fn source_attestation(&self) -> SourceAttestation;
}
\end{lstlisting}

\noindent
where \texttt{SourceAttestation} encapsulates a digital signature
produced by the inference service over the payload, using a key
stored in an HSM or TPM.
\texttt{EvidenceToken::new} would then accept
\texttt{impl AttestableContext} instead of \texttt{\&[u8]}.
Forging evidence would require compromising the signing key---a
problem reducible to PKI security---rather than simply passing
an arbitrary byte string.

This extension tightens the type law to
$V \multimap (E_{\text{attested}} \otimes C)$,
which is strictly stronger than $V \multimap (E \otimes C)$:
every proof of the original invariant is also a proof of the
attested variant, but not vice versa.
The \texttt{Sealed} supertrait prevents external crates from
implementing \texttt{AttestableContext}, preserving the closed
perimeter of Axiom~3.

% ----------------------------------------------------------------
\subsection{Limitations of the Reference Implementation}
\label{sec:limitations}

The \texttt{silent-decisions-proof} kernel is a proof-of-concept
designed to validate the theoretical claims of this paper.
Four implementation choices trade production robustness for
clarity of exposition; each has a well-defined remediation path.

\paragraph{L1 — Static HMAC key (mitigated).}
The reference implementation now reads the HMAC key from the
\texttt{BTV\_HMAC\_KEY} environment variable at runtime, falling back
to a compile-time constant only when the variable is unset.
This closes L1 for deployments that inject a key from an HSM or KMS.
However, L1 also surfaces an architectural asymmetry that recurs
whenever a tamper-evident audit log is deployed alongside BTV: the
verdict integrity seal offers its strongest guarantee only when
\texttt{BTV\_HMAC\_KEY} is provisioned from the same HSM or KMS
infrastructure that protects the audit log's signing key.
Without such alignment, a deployment's overall non-repudiation
guarantee is bounded by its weakest key-management practice.
The type-level guarantee (Theorem~\ref{thm:enclosure}) is independent
of key management and remains valid regardless of the HMAC key source.

\paragraph{L2 — \texttt{ComplianceAuthority} factory (implemented).}
The \texttt{ComplianceToken::new} constructor was originally
\texttt{pub}, allowing any external crate to self-declare a
jurisdiction and policy version without validation.
The reference implementation now restricts this constructor to
\texttt{pub(crate)} and exposes a \texttt{ComplianceAuthority}
factory that validates jurisdiction strings against an allowlist
of recognised regulatory regimes (\texttt{BR-LGPD},
\texttt{EU-GDPR}, \texttt{EU-AI-ACT}).
The authority's signing key is read from the
\texttt{BTV\_AUTHORITY\_KEY} environment variable, with a
deterministic fallback for the proof test suite.
Two additional proof clauses (Clauses~16 and~17) verify that
(a)~an unrecognised jurisdiction (\texttt{"Narnia"}) is rejected,
and (b)~valid jurisdictions produce well-formed tokens.
The type law tightens to
$V \multimap (E \otimes C_{\text{signed}})$; the encapsulation
boundary (Axioms~1--3) is preserved.

\paragraph{L3 — Single-node evidence chain.}
The current \texttt{EvidenceToken} captures a single BLAKE3 hash
of the decision context.
Multi-step AI pipelines (retrieval, re-ranking, generation,
scoring) produce intermediate states that each carry accountability
implications.
Extending the type system to a \emph{chain} of evidence tokens
where each pipeline stage consumes the prior token and produces
a new one is a natural generalisation that preserves
$V \multimap (E \otimes C)$ at each step.

\paragraph{L4 — Polyglot pipelines.}
Production AI systems rarely execute in a single-language process:
Python orchestrators, C++ inference engines, and Java API gateways must
all participate in the same governance boundary.
Because the theorem holds within Safe Rust (Axiom~\ref{ax:encap}),
external code constitutes $\Gamma_{\text{ext}}$ by definition, and
cannot construct a \texttt{Verdict} directly.
Section~\ref{sec:ffiboundary} analyses this boundary as a threat
model: the PyO3 gateway (\texttt{btv-python}) keeps external callers
inside $\Gamma_{\text{ext}}$, while a \emph{BTV sidecar} deployed
alongside each service in a multi-service pipeline can intercept
decision calls and chain a fresh \texttt{EvidenceToken} at each stage,
addressing L3 and L4 jointly.

% ----------------------------------------------------------------
\subsection{CAL Design Trade-Offs}
\label{sec:cal}

A production governance kernel faces a three-way tension among
\textbf{C}onsistency (every materialized verdict carries a
causally-bound evidence chain), \textbf{A}vailability (the system can
always issue a verdict in finite time), and \textbf{L}egality (tokens
are issued only by a verified authority).
Existing audit-logging designs effectively choose
consistency-with-availability over legality: the decision path keeps
issuing decisions even when the logging pipeline fails or drops
records, so the resulting log carries no enforceable evidentiary
guarantee.
BTV makes the opposite choice: when evidence cannot be produced or
the compliance authority is unreachable, the system fails securely
rather than issuing an unevidenced verdict, accepting reduced
availability---a deliberate posture for decisions with legal
consequences, partially recovered by the \texttt{EscalatedVerdict}
path (Section~\ref{sec:escalated}).
We present this as a practical design observation, not a formal
impossibility result; whether an FLP-style impossibility
theorem~\cite{flp1985} holds for accountability systems remains an
open question.

% ----------------------------------------------------------------
\subsection{Future Work}
\label{sec:future}

Six directions are most immediate.
First, implementing the \texttt{AttestableContext} extension
(Section~\ref{sec:scope}) and the \texttt{ComplianceAuthority}
factory (L2) would close the semantic-truthfulness and
self-declaration gaps in the reference kernel.
Second, extending the evidence chain to multi-step pipelines via
the sidecar pattern (L3, L4) would make the framework applicable
to retrieval-augmented generation and agentic AI architectures.
Third, formalising the type-system extension in a proof assistant
(Lean~4 or Coq) would elevate the Constitutional Enclosure
Theorem from a machine-checked Rust proof to a fully mechanised
mathematical proof.
Fourth, investigating whether a formal CAL impossibility result
holds under asynchronous network conditions would provide
theoretical grounding for the design trade-offs observed in
Section~\ref{sec:cal}.
Fifth, unifying \texttt{Verdict} and \texttt{EscalatedVerdict}
under the \texttt{AccountableDecision} trait and formalising the
composed invariant $(V_{\mathrm{auto}} \oplus V_{\mathrm{esc}})
\multimap R$, where $R$ is the linear resource appropriate to
each branch, would complete the type-theoretic account of
mixed human-automated decision pipelines.
Sixth, formalising the composition of BTV's construction-time
guarantee with complementary mechanisms for tamper-evident persistence
and privacy-preserving redaction of transparency logs remains an open
problem.  We conjecture that these mechanisms operate at distinct
enforcement layers (construction-time, delivery-time, and audit-time
respectively) with non-overlapping trust assumptions, making
end-to-end composition plausible; proving this formally---ideally via
a mechanised proof assistant---is left as future work.
